
\chapter{Implementation and Experiments}\label{ch:impl}

This chapter reports the implementation of the design presented in Chapter~\ref{ch:design} and describes the experimental setup used to evaluate it. The purpose of the chapter is to make the artefact concrete: to show how the five-phase pipeline of Section~\ref{sec:design-overview} is realised inside the runtime, how the three mechanism-specific sub-pipelines map to source-tree changes, how the correctness of the scanner is defended through unit and integration tests, and how the end-to-end evaluation is exercised through a dedicated continuous-integration repository built around the runtime.

Section~\ref{sec:impl-overview} gives an overview of the implemented prototype and introduces the call-flow diagram that ties the runtime entry points to the scanner components. Section~\ref{sec:impl-tree} documents the changes the proposed runtime makes to the upstream source tree. Section~\ref{sec:impl-caps} describes the new default capability set and the override flags that keep backwards compatibility. Section~\ref{sec:impl-flags} lists the scanner CLI surface and the bundle layout the scanner produces on disk. Section~\ref{sec:impl-subcmd} details the five sub-commands that carry out the scanner's hooks. Section~\ref{sec:impl-cap-trace} covers the cgroup-wide capability trace and its host-side preconditions. Section~\ref{sec:impl-final} describes the finalisation step. Section~\ref{sec:impl-tests} covers the test strategy, including a dedicated end-to-end repository that runs the scanner on a cloud VM. Section~\ref{sec:impl-setup} specifies the experimental setup. Section~\ref{sec:impl-results} reports the current state of the measurements; Section~\ref{sec:impl-discussion} closes with a short discussion.

\section{Implementation Overview}\label{sec:impl-overview}

The implementation is a fork of the upstream \texttt{runc} reference runtime. The fork keeps the upstream layout intact and adds a small number of new files at the repository root. The entry point of the scanner is plugged into the normal \texttt{runc run} and \texttt{runc create} code paths, so that the scanner is never reached for invocations that do not set the \texttt{-{}-security-scan} flag. When the flag is set, the runtime performs the five-phase pipeline of Section~\ref{sec:design-overview}: in-memory spec relaxation, hook installation, normal container execution, tracer shutdown, and on-disk finalisation.

Figure~\ref{fig:callflow} shows the high-level call flow. The control path through the runtime is unchanged on the upstream side: \texttt{run.go} parses flags, \texttt{create.go} materialises the container, and \texttt{utils\_linux.go} drives the hook lifecycle. The scanner adds three extension points to this flow: a relaxation-and-hook-installation step at spec load time, a set of hook sub-commands that run inside the runtime binary during the container lifecycle, and a finalisation step that runs after the container exits.

\begin{figure}[H]
\centering
\begin{tikzpicture}[
    node distance=0.45cm and 0.55cm,
    entry/.style={rectangle, draw, rounded corners=2pt, minimum width=3.1cm, minimum height=0.85cm, align=center, fill=gray!10, font=\small},
    comp/.style={rectangle, draw, rounded corners=2pt, minimum width=3.1cm, minimum height=0.85cm, align=center, font=\small},
    hook/.style={rectangle, draw, rounded corners=2pt, minimum width=3.3cm, minimum height=0.75cm, align=center, font=\footnotesize},
    artefact/.style={cylinder, shape border rotate=90, aspect=0.25, draw, minimum width=2.6cm, minimum height=0.85cm, align=center, font=\footnotesize},
    arrow/.style={-Stealth, thick},
    dashedarrow/.style={-Stealth, thick, dashed}
]
    \node[entry] (run) {\texttt{runc run}};
    \node[entry, right=of run] (create) {\texttt{runc create}};
    \node[comp, below=0.8cm of $(run)!0.5!(create)$] (scanner) {\texttt{scanner\_linux.go}\\relax spec, install hooks};
    \node[comp, left=1.1cm of scanner] (defaults) {\texttt{utils\_linux.go}\\capability defaults};
    \node[comp, right=1.1cm of scanner] (final) {\texttt{scanner\_finalize\_linux.go}\\narrow \texttt{config.json}};

    \node[hook, below=1.0cm of defaults] (aaload) {\texttt{scan-aa-load}};
    \node[hook, right=0.15cm of aaload] (capstart) {\texttt{scan-cap-trace-start}};
    \node[hook, right=0.15cm of capstart] (snap) {\texttt{scan-cap-snapshot}};
    \node[hook, below=0.25cm of aaload] (capstop) {\texttt{scan-cap-trace-stop}};
    \node[hook, right=0.15cm of capstop] (aaunload) {\texttt{scan-aa-unload}};

    \node[artefact, below=0.75cm of capstop] (log) {\texttt{capable-}\\\texttt{bpfcc.log}};
    \node[artefact, right=0.25cm of log] (sec) {\texttt{seccomp.}\\\texttt{json}};
    \node[artefact, right=0.25cm of sec] (aa) {\texttt{apparmor.}\\\texttt{profile}};

    \draw[arrow] (run) -- (scanner);
    \draw[arrow] (create) -- (scanner);
    \draw[arrow] (defaults) -- (scanner);
    \draw[arrow] (scanner) -- (final);
    \draw[dashedarrow] (scanner.south) -- ++(0,-0.25) -| (aaload.north);
    \draw[dashedarrow] (scanner.south) -- ++(0,-0.25) -| (capstart.north);
    \draw[dashedarrow] (scanner.south) -- ++(0,-0.25) -| (snap.north);
    \draw[dashedarrow] (scanner.south) -- ++(0,-0.25) -| (capstop.north);
    \draw[dashedarrow] (scanner.south) -- ++(0,-0.25) -| (aaunload.north);
    \draw[arrow] (capstart) -- (log);
    \draw[arrow] (capstop) -- (log);
    \draw[arrow] (scanner) -- (sec);
    \draw[arrow] (aaload) -- (aa);
    \draw[arrow] (log) -- ++(0,-0.25) -| (final.south);
\end{tikzpicture}
\caption{Call-flow view of the proposed runtime. Solid arrows show direct function calls; dashed arrows show OCI hooks whose \texttt{Path} points back at the runtime binary and whose argument names the sub-command to execute.}
\label{fig:callflow}
\end{figure}

The same figure makes the role of each hook concrete: the AppArmor sub-pipeline runs between \texttt{scan-aa-load} (createRuntime) and \texttt{scan-aa-unload} (poststop), the capability sub-pipeline runs between \texttt{scan-cap-trace-start} (poststart) and \texttt{scan-cap-trace-stop} (poststop, prepended), and the snapshot sub-command runs once after the container starts to capture \texttt{/proc/<init>/status} for diagnostics. The seccomp sub-pipeline is invoked through an upstream prestart hook that the scanner installs; its output is written by the hook itself, without a dedicated sub-command.

End-to-end testing is not run from inside the runtime repository. A dedicated companion repository, \texttt{thesis-ci-repo}, contains the matrix workflow that exercises the scanner against a growing list of bundles on a cloud VM; the repository packages the cloud infrastructure, the host provisioning, the per-bundle contract, and the artefact archive as code. Its contents are documented as an additional achievement in Chapter~\ref{ch:additional}; here, the companion repository plays the role an \emph{end-to-end test suite} plays for a standalone application: it checks that the scanner behaves correctly against real bundles, real tooling, and real kernels rather than only against the mocks the in-repository integration tests use.

\section{Source-Tree Changes}\label{sec:impl-tree}

The new code is organised into four Go files at the repository root, with table-driven unit tests next to each. The files and their responsibilities are listed in Table~\ref{tab:sourcechanges}.

\begin{table}[ht]
\centering
\caption{Files added or modified by the proposed runtime.}
\label{tab:sourcechanges}
\begin{tabular}{p{4.8cm}p{9.4cm}}
\toprule
\textbf{File} & \textbf{Responsibility} \\
\midrule
\texttt{scanner\_linux.go} & Entry point of the scanning mode: loads the spec, relaxes it in memory, grants all capabilities, installs hooks, calls into the finalisation step after the container exits. \\
\texttt{scanner\_hooks\_linux.go} & Implementation of the five sub-commands (\texttt{scan-aa-load}, \texttt{scan-aa-unload}, \texttt{scan-cap-snapshot}, \texttt{scan-cap-trace-start}, \texttt{scan-cap-trace-stop}). \\
\texttt{scanner\_bpf\_linux.go} & cgroup-v2 resolution, BPF-map creation and pinning through \texttt{bpftool}, feature detection for \texttt{capable-bpfcc}. \\
\texttt{scanner\_finalize\_linux.go} & Parser for the capability log and the narrowing step on the on-disk \texttt{config.json}. \\
\texttt{scanner\_*\_test.go} & Table-driven unit tests for each of the four files above. \\
\texttt{run.go}, \texttt{create.go} & CLI flag registration (\texttt{-{}-security-scan}, \texttt{-{}-default-capabilities}, \texttt{-{}-default-capabilities-docker}, \texttt{-{}-scan-capable}, \texttt{-{}-scan-bpftool}, \texttt{-{}-scan-seccomp-hook}). \\
\texttt{libcontainer/specconv/example.go} & New default capability template (empty set). \\
\texttt{utils\_linux.go} & Override logic for capability defaults and the entry point into the finalisation step. \\
\bottomrule
\end{tabular}
\end{table}

The upstream \texttt{runc} source-tree layout is otherwise left alone. The scanner files are compiled only on Linux (\texttt{\_linux.go} suffix) so that non-Linux builds of \texttt{runc} still compile without the scanner. Flag registration in \texttt{run.go} and \texttt{create.go} is additive: if the operator does not pass any of the scanner flags, the only code the runtime executes at spec-load time is the new capability-default helper, which reduces to a single \texttt{switch} on a string that is empty in the unflagged case.

\section{Capability Defaults and Override Logic}\label{sec:impl-caps}

The default \texttt{runc spec} template now emits an empty \texttt{capabilities} object. The behavioural change is visible from the command line:

\begin{verbatim}
$ runc spec
$ jq '.process.capabilities' config.json
{}
\end{verbatim}

Two flags restore the historical defaults. They are mutually exclusive:

\begin{verbatim}
runc run --default-capabilities          <id>   # 3 caps  (upstream)
runc run --default-capabilities-docker   <id>   # 14 caps (Docker)
runc run                                 <id>   # 0 caps  (new default)
\end{verbatim}

The override logic lives in \texttt{utils\_linux.go}. A helper returns one of three modes (empty, \texttt{runc}, \texttt{docker}) from the CLI context; passing both flags returns an error from the CLI parser rather than from the container start path, so the operator sees the problem before the spec is materialised. The override rewrites all five capability buckets (bounding, effective, permitted, inheritable, ambient) so that capabilities survive the \texttt{exec(2)} boundary consistently; writing only the bounding set would cause the effective set to lose capabilities on re-exec.

The resulting capability masks for the three modes are summarised in Table~\ref{tab:capmodes}. The masks were recorded by decoding the \texttt{CapPrm} field of \texttt{/proc/<init>/status} after starting a container in each mode.

\begin{table}[ht]
\centering
\caption{Effective capability mask for the three default modes.}
\label{tab:capmodes}
\begin{tabular}{p{4.8cm}p{3.6cm}p{5.2cm}}
\toprule
\textbf{Mode} & \textbf{CapPrm} & \textbf{Interpretation} \\
\midrule
Default (no flag) & \texttt{0x0000000000000000} & empty \\
\texttt{-{}-default-capabilities} & \texttt{0x0000000020000420} & \texttt{CAP\_KILL}, \texttt{CAP\_NET\_BIND\_SERVICE}, \texttt{CAP\_AUDIT\_WRITE} \\
\texttt{-{}-default-capabilities-docker} & \texttt{0x00000000a80425fb} & historical Docker fourteen-cap set \\
Both flags & --- & CLI returns ``mutually exclusive'' error \\
\bottomrule
\end{tabular}
\end{table}

\section{Scanner CLI and Bundle Layout}\label{sec:impl-flags}

The scanner is opt-in. The entry point is a single flag; the remaining flags are override paths for the external tooling, useful on hosts that install \texttt{capable-bpfcc} or \texttt{bpftool} outside \texttt{PATH}.

\begin{verbatim}
runc run \
  --security-scan \
  [--scan-seccomp-hook PATH] \
  [--scan-capable PATH] \
  [--scan-bpftool PATH] \
  <container-id>
\end{verbatim}

When \texttt{-{}-security-scan} is set, the bundle gains a single new directory next to the rootfs. No executable is ever written into the bundle; the scanner's own hooks are sub-commands of the runtime binary and the external tracing tools are invoked from inside those sub-commands. The produced tree is shown below:

\begin{verbatim}
<bundle>/generated/
  seccomp.json                          # allow-list profile
  capable-bpfcc.log                     # raw cgroup-wide cap trace
  capabilities-from-proc-status.txt     # poststart /proc/<init>/status snapshot
  apparmor.profile                      # complain-mode profile
  apparmor-load.log                     # apparmor_parser stdout/stderr
  apparmor-README.txt                   # how to refine via aa-logprof
  spec.original.json                    # pre-scan config.json, kept for rollback
\end{verbatim}

The bundle layout is stable across scanner versions and is the contract that the end-to-end tests in the companion repository check against.

\section{Self-Exec Hook Sub-Commands}\label{sec:impl-subcmd}

Five hidden sub-commands are registered next to \texttt{run}, \texttt{create}, and \texttt{exec} in the CLI. Each accepts the OCI hook state JSON on stdin and reads its parameters from environment variables propagated through the OCI \texttt{Hook.Env} field. The sub-commands are summarised in Table~\ref{tab:subcmds}.

\begin{table}[ht]
\centering
\caption{Sub-commands registered by the proposed runtime.}
\label{tab:subcmds}
\begin{tabular}{p{3.8cm}p{2.5cm}p{7.8cm}}
\toprule
\textbf{Sub-command} & \textbf{Hook stage} & \textbf{Responsibility} \\
\midrule
\texttt{scan-aa-load} & \texttt{createRuntime} (prepended) & Calls \texttt{apparmor\_parser -r} on the generated profile; saves stdout/stderr to \texttt{generated/apparmor-load.log}. \\
\texttt{scan-aa-unload} & \texttt{poststop} & Calls \texttt{apparmor\_parser -R} to remove the profile. \\
\texttt{scan-cap-snapshot} & \texttt{poststart} & Reads \texttt{/proc/<init>/status}, saves the \texttt{Cap*} lines into \texttt{capabilities-from-proc-status.txt} for diagnostics. \\
\texttt{scan-cap-trace-start} & \texttt{poststart} & Resolves the cgroup path and inode, creates a hash map under \texttt{/sys/fs/bpf/runc-scan/<id>/cgroup\_filter} via \texttt{bpftool map create}, fills it via \texttt{bpftool map update}, \texttt{Setpgid}-spawns \texttt{capable-bpfcc -{}-cgroupmap <map>}. Writes the tracer PID into \texttt{generated/.runc\_cap\_trace.pid}. \\
\texttt{scan-cap-trace-stop} & \texttt{poststop} (prepended) & \texttt{SIGTERM}s the tracer, waits up to two seconds for the log to flush, escalates to \texttt{SIGKILL} if needed, removes the pin under \texttt{/sys/fs/bpf/runc-scan/<id>/}. \\
\bottomrule
\end{tabular}
\end{table}

Three implementation points are worth recording. First, every sub-command drains stdin before returning, because the OCI runtime pipes the hook state JSON even to hooks that do not consume it, and failing to read it risks a \texttt{SIGPIPE} on the runtime side. Second, every sub-command writes its own log file under \texttt{generated/} so that a failure in one hook is diagnosable without the other hooks' output being interleaved in the same log. Third, the \texttt{scan-cap-trace-stop} sub-command is prepended to the poststop list so that it runs before any user-defined hook or the AppArmor unload hook; this ordering makes sure the in-kernel ring buffer is flushed before any hook that might disturb the cgroup state.

\section{cgroup-Wide Capability Trace}\label{sec:impl-cap-trace}

A host-side preconditions check runs before the scanner actually starts; it decides whether the scanner can run at all. The check inspects five invariants, listed in Table~\ref{tab:preconditions}.

\begin{table}[ht]
\centering
\caption{Host-side preconditions for the capability trace.}
\label{tab:preconditions}
\begin{tabular}{p{5.0cm}p{9.2cm}}
\toprule
\textbf{Precondition} & \textbf{Verification} \\
\midrule
\texttt{capable-bpfcc} available & On \texttt{PATH} or supplied via \texttt{-{}-scan-capable}. \\
\texttt{capable-bpfcc -{}-cgroupmap} supported & Verified by parsing \texttt{capable-bpfcc -{}-help} output. \\
\texttt{bpftool} available & On \texttt{PATH} or supplied via \texttt{-{}-scan-bpftool}. \\
Unified cgroup-v2 hierarchy & Verified by reading \texttt{/proc/mounts}; hybrid hosts are rejected. \\
\texttt{/sys/fs/bpf} is a \texttt{bpffs} mount & Verified by reading \texttt{/proc/mounts}; skipped under a test-only environment variable that redirects the pin root, so that CI can run under a mocked bpffs. \\
\bottomrule
\end{tabular}
\end{table}

If any precondition fails, scan mode refuses to run and returns an actionable error that names the missing tool or the failing check. There is no silent fallback to the PID-based tracing mode of \texttt{capable-bpfcc}, because that fallback would degrade child-process coverage and silently break requirement R2 from Chapter~\ref{ch:design}.

\section{Finalisation}\label{sec:impl-final}

The finalisation function is called from \texttt{utils\_linux.go} after the container exits. Its four responsibilities are:

\begin{enumerate}
    \item parse \texttt{generated/capable-bpfcc.log}, extracting capability names from the kernel-emitted \texttt{cap\_capable} events;
    \item normalise each observed name through a single helper that maps BCC's lower-case short names to OCI's upper-case \texttt{CAP\_*} form;
    \item replace \texttt{Process.Capabilities.\{Bounding, Effective, Permitted, Inheritable, Ambient\}} on the \emph{on-disk} \texttt{config.json} with the observed set;
    \item write the previous \texttt{config.json} verbatim to \texttt{generated/spec.original.json}, so that a later run can roll back the narrowing by copying the backup over the active spec.
\end{enumerate}

The function is set-difference, not set-union: an empty trace produces an empty capability set, with an informational log line that names the log file that produced the empty result. This is the safe reading of the evidence, and it matches requirement R5 in Chapter~\ref{ch:design}: re-running the scanner never widens the eventual enforcement spec.

\section{Tests}\label{sec:impl-tests}

The correctness of the scanner is defended at three levels: unit tests inside the runtime repository, an integration smoke pack inside the runtime repository, and an end-to-end matrix inside the companion CI repository. Table~\ref{tab:tests} summarises the three levels.

\begin{table}[ht]
\centering
\caption{Test suites that exercise the proposed runtime.}
\label{tab:tests}
\begin{tabular}{p{3.6cm}p{4.0cm}p{7.6cm}}
\toprule
\textbf{Suite} & \textbf{Scope} & \textbf{What it checks} \\
\midrule
Unit tests & In-process Go tests & Capability-name normaliser, cgroup-path resolver, AppArmor profile-name sanitiser, finalisation parser, relax-spec mutator, grant-all-caps mutator. \\
Integration smoke & bats + runtime repo & The CLI flag end-to-end against a busybox bundle; stubs replace \texttt{oci-seccomp-bpf-hook}, \texttt{capable-bpfcc}, and \texttt{bpftool}. \\
End-to-end matrix & \texttt{thesis-ci-repo} & Real bundles, real tooling, real kernel; runs on a self-hosted runner provisioned from infrastructure-as-code. \\
\bottomrule
\end{tabular}
\end{table}

\textbf{Unit tests.} The unit tests are table-driven and live next to the implementation. They cover the pure functions: the capability-name normaliser (forty-odd cases, including the tricky \texttt{CAP\_BPF} vs \texttt{cap\_bpf} distinction), the cgroup-path resolver on both cgroup-v2 layouts produced by systemd-managed and non-systemd hosts, the AppArmor profile-name sanitiser for non-ASCII container identifiers, the finalisation parser on crafted \texttt{capable-bpfcc} log fragments, and the two spec mutators that perform relaxation and grant-all-capabilities. The unit tests run under \texttt{go test ./...} and do not require any external tooling.

\textbf{Integration smoke (in-repo).} A \texttt{bats} suite exercises the CLI flag end-to-end with a minimal \texttt{busybox} bundle and a set of stubs that pretend to be \texttt{oci-seccomp-bpf-hook}, \texttt{capable-bpfcc}, and \texttt{bpftool}. The stubs ship under \texttt{contrib/} and are used only by the smoke pack. The test asserts that \texttt{generated/seccomp.json} and \texttt{generated/capable-bpfcc.log} are created, that \texttt{generated/apparmor.profile} contains the expected \texttt{flags=(complain, \ldots)} marker, that the bundle contains no executable file with the \texttt{.sh} extension (a direct check of requirement R3), and that the on-disk \texttt{config.json} has been narrowed or, for the empty stub trace, set to an empty capability set. The smoke pack is invoked as:

\begin{verbatim}
sudo make localintegration TESTPATH=/security_scan.bats
\end{verbatim}

\textbf{End-to-end matrix (companion repository).} The end-to-end tests live in a separate repository, \texttt{thesis-ci-repo}, documented in Chapter~\ref{ch:additional}. The separation is intentional: the matrix is not part of the runtime under test, and mixing the observer with the observed makes it harder to argue that the matrix is fair. The matrix uses the real \texttt{oci-seccomp-bpf-hook}, the real \texttt{capable-bpfcc} from the BCC toolchain~\cite{bcc2025}, and the real \texttt{bpftool} from \texttt{linux-tools}, running on a cloud VM whose image is managed by Ansible and whose lifecycle is managed by Terraform. The list of bundles grows as new workloads are added; the current catalogue includes a smoke-pack \texttt{busybox-default} and an application-layer \texttt{juice-shop} bundle that runs a live HTTP probe during the scan. Each matrix run keeps the entire \texttt{generated/} directory in object storage, partitioned by runtime commit, run identifier, and bundle name, so that any figure reported in this thesis can be traced back to the exact \texttt{config.json} and the exact log that produced it.

The companion repository plays the role of the end-to-end test harness. It triggers on every push to the runtime fork (through a \texttt{repository\_dispatch} event that carries the runtime commit identifier), on a nightly cron, and on demand from a \texttt{workflow\_dispatch} button. The per-bundle \texttt{verify-profiles.sh} script, plus a bundle-defined \texttt{bats} pack, checks the same contract that the in-repository smoke test checks, but against real kernel behaviour rather than against stubs. A regression in the scanner that passes the in-repository smoke pack (because the stubs cannot tell a real \texttt{oci-seccomp-bpf-hook} apart from their own output) is therefore caught by the companion repository on the next push.

\section{Experimental Setup}\label{sec:impl-setup}

The experimental setup instantiates the evaluation methodology of Section~\ref{sec:design-eval}. The matrix workflow runs on a single self-hosted runner provisioned by the companion repository. The runner is a Yandex Cloud VM with four virtual CPUs and eight gigabytes of memory, running Ubuntu~24.04 with a 6.8-series kernel; the exact machine image is defined in Terraform and reprovisioning it produces a byte-identical host configuration.

Every bundle in the matrix is run under each of the five configurations from Section~\ref{sec:design-eval}: upstream \texttt{runc} in its default state, the proposed runtime with the empty-capability default but no scanning, the proposed runtime after a full scan-then-enforce cycle, gVisor, and Kata Containers. The measurement plan calls for ten repetitions after one warm-up iteration per configuration and per bundle, with the median and the interquartile range reported; wall-clock times are taken from \texttt{runc events} and from bundle-defined probe scripts, CPU and memory are taken from the container cgroup (\texttt{cpu.stat}, \texttt{memory.current}), I/O bandwidth is taken from a synthetic \texttt{fio} job or from a bundle-defined probe, and TCP throughput is taken from an \texttt{iperf3} pair. The exact reproducer commands live in each bundle's \texttt{bundle.yaml} contract in the companion repository, so that any future researcher can replay the matrix without access to the original workstation.

\section{Current Results}\label{sec:impl-results}

The comparative evaluation has not been executed at the time of writing; what follows is the schema under which the numbers will be reported, together with placeholders that mark the tables whose contents depend on the pending runs. The approach keeps the schema visible now so that the reader can tell which rows are instrumental and which rows are expected to be filled in.

\begin{table}[ht]
\centering
\caption{Planned reporting: steady-state CPU overhead relative to default \texttt{runc} ($N=10$).}
\label{tab:cpu}
\begin{tabular}{lrrrr}
\toprule
\textbf{Workload} & \textbf{Default \texttt{runc}} & \textbf{Proposed runtime} & \textbf{gVisor} & \textbf{Kata} \\
\midrule
nginx       & ---  & --- & --- & --- \\
redis       & ---  & --- & --- & --- \\
postgresql  & ---  & --- & --- & --- \\
juice-shop  & ---  & --- & --- & --- \\
\bottomrule
\end{tabular}
\end{table}

\begin{table}[ht]
\centering
\caption{Planned reporting: cold-start latency, ms (median, IQR), $N=10$.}
\label{tab:startup}
\begin{tabular}{lrrrr}
\toprule
\textbf{Workload} & \textbf{Default \texttt{runc}} & \textbf{Proposed runtime} & \textbf{gVisor} & \textbf{Kata} \\
\midrule
busybox     & ---  & --- & --- & --- \\
nginx       & ---  & --- & --- & --- \\
redis       & ---  & --- & --- & --- \\
juice-shop  & ---  & --- & --- & --- \\
\bottomrule
\end{tabular}
\end{table}

\begin{table}[ht]
\centering
\caption{Planned reporting: profile quality emitted by the scanner.}
\label{tab:profquality}
\begin{tabular}{lrrr}
\toprule
\textbf{Workload} & \textbf{Allowed syscalls} & \textbf{Capabilities used} & \textbf{AppArmor entries} \\
\midrule
busybox     & --- & --- & --- \\
nginx       & --- & --- & --- \\
redis       & --- & --- & --- \\
postgresql  & --- & --- & --- \\
juice-shop  & --- & --- & --- \\
\bottomrule
\end{tabular}
\end{table}

\begin{table}[ht]
\centering
\caption{Planned reporting: security effectiveness against curated container-escape attempts (blocked / total).}
\label{tab:security}
\begin{tabular}{lcccc}
\toprule
\textbf{Attempt class} & \textbf{Default \texttt{runc}} & \textbf{Proposed runtime} & \textbf{gVisor} & \textbf{Kata} \\
\midrule
\texttt{CVE-2019-5736} family~\cite{cve2019} & --- & --- & --- & --- \\
Capability-bound escapes         & --- & --- & --- & --- \\
\texttt{ptrace}/\texttt{kexec}-class syscalls & --- & --- & --- & --- \\
Filesystem confinement breakouts & --- & --- & --- & --- \\
\bottomrule
\end{tabular}
\end{table}

The list of vulnerable-container workloads that will be run against the security-effectiveness table is being drawn from the application-security training ecosystem (OWASP Juice Shop, a deliberately misconfigured database container, and comparable targets); the exact aggregation method is still being developed, as noted in Section~\ref{sec:design-eval}.

\section{Discussion}\label{sec:impl-discussion}

The implementation delivers the artefacts the design called for: a single CLI flag, three sub-pipelines that produce four complementary artefacts, a self-exec hook mechanism that keeps the bundle free of executable files, and a finalisation step whose semantics is set difference. The five preconditions in Section~\ref{sec:impl-cap-trace} make concrete the trade-off the design chapter stated explicitly --- prefer a clean refusal to a silent degradation --- and are the single most important safety invariant in the implementation. Once the comparative matrix completes, the qualitative claim that the hardened runtime closes the kernel-side perimeter against the syscall-, capability-, and path-based attack classes while staying within single-digit percent of default \texttt{runc} for CPU and memory can be checked quantitatively; the discussion of the trade-offs, the sensitivity of the numbers to workload choice, and the comparison to the sandboxed alternatives is deferred to Chapter~\ref{ch:conclusion}.
